% Evidence-safe revision of the July 2026 spectral/entropy manuscript.
% This file deliberately excludes all unverified end-to-end benchmark claims.
% RETARGET (2026-10-19 panel, EIC-W1): negative-result / reproducibility
% workshop mechanism audit (prototype autopsy), NOT a main-track novelty
% submission. Claims are shrunk to existing evidence; no new training runs.
% Canonical copy: this file. The divergent copy at
% zvf-program/flagship/paper/spectral_entropy_paper_kdense_revision.tex is
% SUPERSEDED (see header there); do not submit it.
% AUTHOR_TODO (EIC-W3): once the workshop venue is bound, produce the
% anonymized variant (with neurips_2026.sty this is dropping the `preprint'
% option, which switches to anonymous submission mode), vendor the style
% file or state build instructions for the bound venue, and add its checklist.
\ifdefined\pdfsuppressptexinfo\pdfsuppressptexinfo=-1\fi
\ifdefined\pdfinfoomitdate\pdfinfoomitdate=1\fi
\ifdefined\pdftrailerid\pdftrailerid{}\fi

\documentclass{article}

\usepackage[preprint,nonatbib]{neurips_2026}
\usepackage[numbers,sort&compress]{natbib}
\usepackage[utf8]{inputenc}
\usepackage[T1]{fontenc}
\usepackage{amsmath,amssymb,amsfonts,amsthm}
\usepackage{booktabs}
\usepackage{enumitem}
\usepackage{microtype}
\usepackage{tabularx}
\usepackage{xcolor}
\usepackage{xspace}
\usepackage{hyperref}
\usepackage{url}

\hypersetup{
  colorlinks=true,
  linkcolor=blue!55!black,
  citecolor=blue!55!black,
  urlcolor=blue!55!black
}

\newtheorem{proposition}{Proposition}
\newtheorem{lemma}{Lemma}
\newtheorem{remark}{Remark}
\newcommand{\grpo}{\textsc{GRPO}\xspace}
\newcommand{\zvf}{\textsc{ZVF}\xspace}
\newcommand{\E}{\mathbb{E}}
\newcommand{\Prob}{\mathbb{P}}

\title{Zero Advantage Is Not Missing Information:\\
A Mechanism Audit of Trajectory-Spectral Auxiliary Advantages}

\author{%
  Arvind C R \\
  PES University \\
  \texttt{arvindcr4@gmail.com}
}

\begin{document}
\maketitle

\begin{abstract}
Group-relative policy optimization assigns zero centered advantage whenever
every completion in a sampled group receives the same reward.  We argue that
this is an absence of task feedback rather than hidden information.  No
transform placed only in the forward path can restore a policy gradient once
the centered advantage vanishes; only an explicit auxiliary advantage can,
and that changes the objective.  We dissect one such prototype: a standardized
score built from pairwise distances between low-order Legendre features of
sampled trajectories, with an adaptive Givens rotation for coordinate
elimination.  This is a negative-result mechanism audit (a prototype
autopsy with transferable checks), not a main-track novelty claim. The
continuous Parseval identity holds only for the reconstructed truncated
expansion, while the implemented equispaced rule is first-order.  The
rotation preserves norm precisely because its subsequent projection discards
an already annihilated coordinate in exact arithmetic. The
$2^{1-G}$ homogeneous-group floor is a one-line corollary of the binary
identity and Jensen bound in \citep{nie2026gradient}, not a new result
here. On synthetic fixtures the coupled objective produces a gradient
proportional to its coupling weight, but the fixture injects the
between-completion offsets the score recovers and its position channel
cancels in every pairwise distance, so this $\lambda$-linearity is an
algebraic identity check, not validation. The masked/padded defect of
Eq.~\eqref{eq:discrete-coeff} is unmeasured by the all-active replay
below. The entropy gate as implemented never fires (0/18 comparisons)
because its threshold $\tau=0.5$ is an order of magnitude below the
$\log L$-scale statistic it gates; this is a scaling defect of the
implementation, not evidence about entropic gating as a concept. We
specify prospective alignment, placebo, paired-seed, and ablation tests
and report no benchmark result.
\end{abstract}

\section{Introduction}

Group Relative Policy Optimization and related critic-free objectives compare
several completions for the same prompt and standardize their rewards within
the sampled group \citep{shao2024deepseekmath}.  If all rewards agree, the
centered reward term is exactly zero.  Recent work studies this phenomenon as
gradient starvation or advantage collapse and proposes changes to sampling,
normalization, or the objective \citep{nie2026gradient,he2026advantage}.

An equal-reward group is observationally ambiguous.  All completions may be
wrong, all may be correct, the verifier may be coarse or mistaken, or the
completions may differ in ways irrelevant to task success.  A representation
distance can reveal that trajectories differ, but cannot by itself reveal
which difference is useful.  Consequently, adding such a distance to the
advantage is a new auxiliary objective.  It is not recovery of a missing
reward, and a nonzero gradient is not evidence that the gradient is beneficial.

This paper evaluates a narrow prototype under that interpretation. It is
written as a negative-result mechanism audit --- an autopsy of the authors'
own prototype whose transferable content is the checks, not the score ---
for a negative-results/reproducibility venue, and makes no main-track
novelty claim. It makes five evidence-bounded audit findings:

\begin{enumerate}[leftmargin=*,itemsep=2pt]
  \item We separate a missing reward signal from an auxiliary objective: no
  forward-path transform can create an update when its centered advantage is
  zero, whereas explicit auxiliary coupling changes the optimized objective.
  \item We define a trajectory-spectral auxiliary score and locate the
  exactness boundary between continuous Legendre identities and the implemented
  equispaced discretization.
  \item We audit adaptive Givens elimination and show that its projection is
  redundant by construction: the rotation already annihilates the coordinate
  it projects away.
  \item We restate, as a credited corollary (see \S3), the conditionally
  i.i.d. binary-reward feasibility check
  $\Prob(\text{homogeneous})\geq 2^{1-G}$ and state what procedural changes are
  required for any lower reported rate. We claim no novelty for this floor.
  \item We provide reproducible synthetic identity checks (not validation:
  see \S6) and specify prospective alignment, placebo, paired-seed,
  equal-compute, and ablation gates. We make no policy-learning claim.
\end{enumerate}

\paragraph{Evidence status, stated before any result.}
This paper contains proofs, executable objective modules, 16 focused
unit/harness tests rerun on CPU, and a new deterministic replay of those
modules. The historical independent audit bundle is unavailable in this
checkout; its reported case count and numerical quadrature, norm-loss, and
gradient results are withdrawn rather than treated as verified evidence.
The replacement replay has nine fixture cells and is not a reconstruction
of that missing bundle. Source hashes, commands, and result fields are
listed in Appendix~\ref{app:artifact}.
There is no language-model fine-tuning, held-out benchmark accuracy,
paired-seed learning comparison, or calibration of the trajectory score against
an external target. We make no benchmark-accuracy, sample-efficiency,
deployment, mitigation, or superiority claim.

\paragraph{Companion-manuscript disclosure.}
A companion manuscript, ``Same Terminal Signal, Different Action''
(\path{zvf-program/flagship/paper/main.tex}), shares with this audit the
zero-advantage ambiguity framing, the binary homogeneous-group identity,
and the same \path{pilot/objective.py} artifact. Shared: the ambiguity
thesis and the objective artifact. Unique to the companion: the
fail-closed conformance chain and its screening campaign. Unique here:
the Legendre/Givens dissection, the synthetic replay, and the prospective
alignment gate for the spectral score. Readers and reviewers of either
text should not credit the shared framing twice. This file is the
canonical copy of this audit; the divergent copy under
\path{zvf-program/flagship/paper/} is marked superseded and must not be
submitted, as it retains withdrawn language.

\section{Related Work and Scope}

\grpo removes a learned critic by constructing within-group advantages
\citep{shao2024deepseekmath}.  DAPO combines dynamic sampling with a
token-level loss design \citep{yu2025dapo}; recent analyses isolate
group-variance pathologies and propose direct remedies
\citep{nie2026gradient,he2026advantage}.  These methods motivate the problem
but do not validate the trajectory-derived auxiliary score studied here.

Orthogonal polynomial expansions are classical tools for representing
functions on compact domains.  HiPPO and Legendre Memory Units use scaled
Legendre projections for sequence memory, and later work studies the discrete
basis and parallel training choices directly
\citep{gu2020hippo,voelker2019lmu,chilkuri2021parallelizing}.  PoPE applies
Legendre polynomials to the token axis of a language model
\citep{aggarwal2024pope}.  The closest work by purpose uses within-group
hidden-state distances for span-level credit assignment and, critically,
specifies how that signal enters the advantage \citep{chen2026shear}.
Entropy-guided advantage shaping is also an existing response to zero-variance
prompts \citep{le2025noprompt}.  Our use is therefore a narrower recombination:
a masked trajectory projection supplies a group-relative auxiliary score.  We
do not claim that low modes are semantic or high modes are noise without an
external validation target.

Norm preservation through composed Givens rotations is established
\citep{dorobantu2016dizzyrnn,jing2017eunn,frerix2019givens}.  Planar rotations
already occur in transformer attention through rotary position embedding
\citep{su2024roformer}, and have been used for parameter-efficient adaptation,
compression, and transformer blocks
\citep{ma2024qoft,su2025rotatekv,abuayyash2026threephase}.  Attention-entropy
collapse is likewise a known transformer failure mode \citep{zhai2023entropy}.
The implementation studied here uses a data-dependent rotation as a
coordinate-elimination primitive, not a quantum algorithm. The older source
retained below the disabled entry point in \path{spectral_entropy_paper.tex}
used a ``quantum-inspired'' label that this revision withdraws. We claim no
novelty for its constituents. The contribution is a negative mechanism audit
plus a proposed trajectory-spectral coupling whose usefulness remains to be tested.
All motivation below is scoped to the two-channel log-probability fixture
actually replayed; the supplied-hidden-states path is implemented but
entirely unexercised by any number in this paper, so no hidden-state,
layer-choice, or large-$D$ claim is made. No published external system is
shown here to commit the audited confusions; the audited artifact is the
authors' own prototype, and the paper should be read as a prototype
autopsy whose floor check and claim-index practice are the transferable
parts.

\section{Problem Statement and Causal Boundary}

For prompt $x$, let $y_1,\ldots,y_G$ be sampled completions with scalar rewards
$r_1,\ldots,r_G$. Let $I$ be the selected completion rows, $a=|I|\ge2$,
and require positive masked length for every selected row. All means and
standard deviations below use only $I$:
\[
 \bar u=a^{-1}\sum_{i\in I}u_i,\qquad
 \sigma_u=\sqrt{\frac{1}{a-1}\sum_{i\in I}(u_i-\bar u)^2}.
\]
The current synthetic objective defines, for $i\in I$,
\begin{equation}
  A_i^{\mathrm{reward}}
  = \begin{cases}(r_i-\bar r)/\sigma_r,&\sigma_r>10^{-12},\\
  0,&\text{otherwise}.\end{cases}
\end{equation}
Unselected advantages are zero. Unlike the auxiliary below, this reward
denominator has no additive epsilon. Single-selected-row sample variance is
undefined and is outside the replay's contract. The code contract is
narrower than the notation above: for $G=8$ the artifact accepts exactly
6 or 8 active rows and rejects other selections, and it permits a single
active row, for which the sample standard deviation is NaN and NaN
advantages can propagate instead of the zero stated here. The $a\ge2$
precondition is therefore disclosed but not enforced by the artifact;
enforcing it with an explicit error is a required repair before any
training use, and no result below relies on $a<2$ or on selections the
artifact rejects.

\begin{remark}[Homogeneous rewards remove the centered reward signal]
If $r_1=\cdots=r_G$, then $A_i^{\mathrm{reward}}=0$ for every row, since
homogeneity gives $r_i=\bar r$ and every numerator is zero. Any
score-function term weighted only by these advantages is therefore zero.
This is arithmetic, stated only to fix terminology.
\end{remark}

For binary rewards with per-completion success probability $p$, the probability
that an independent group is homogeneous is
\begin{equation}
  \Prob(\text{homogeneous}\mid p,G)=p^G+(1-p)^G.
\end{equation}
This identity diagnoses when reward contrast is unavailable.  It neither
identifies the latent cause nor supplies a replacement preference.
For binary rewards, the smallest possible homogeneous-group probability over
$p$ is $2^{1-G}$, attained at $p=1/2$.  In particular, $G=4$ has a $12.5\%$
floor.  Any lower reported rate must therefore use a larger group, non-binary
rewards, non-i.i.d. joint sampling, dynamic rejection/resampling, or a different denominator; it cannot
be compared to this identity without stating that procedural change.
This floor is a one-line corollary of the identity above and of the
stronger Jensen lower bound proved in \citep{nie2026gradient}, which also
reports empirical degeneracy rates from training logs; we claim no novelty
for it. Prompt-mixture caveat: with prompt-dependent success probability
$p_x$, the operational homogeneous rate is
$\E_x[p_x^G+(1-p_x)^G]$, governed by the difficulty distribution and
generally well above the pointwise floor. The floor can therefore falsify
only a reported rate below $2^{1-G}$ under the stated i.i.d.\ binary
procedure; it cannot predict, explain, or bound tightly the rate of any
real prompt mixture, and a rate above the floor is uninformative about
which remedy, if any, is appropriate.

The current \path{pilot/objective.py} defines a trajectory score $s_i$ and uses
\begin{equation}
  A_i^{\mathrm{aux}}
  = \begin{cases}(s_i-\bar s)/(\sigma_s+10^{-4}),&\sigma_s>10^{-12},\\
  0,&\text{otherwise},\end{cases}
  \qquad
  \widetilde A_i
  = \operatorname{stopgrad}\!\left(A_i^{\mathrm{reward}}+\lambda A_i^{\mathrm{aux}}\right).
  \label{eq:combined-advantage}
\end{equation}
Equation~\eqref{eq:combined-advantage} is an objective modification.  When
rewards are homogeneous, it can create a nonzero update if the trajectory
scores differ, but its direction is justified only if $s_i$ is shown to align
with future task reward or a valid process-level target.
Detachment occurs after standardization and combination, not before score
construction. Write $z_{i\ell}$ for the input log-probability tensor and
$\rho_{i\ell}=\exp(z_{i\ell}-z^{\rm old}_{i\ell})$. The clipped
surrogate follows the usual PPO form \citep{schulman2017ppo}. The minimized loss is
\begin{equation}
 \mathcal L=-\frac1a\sum_{i\in I}\frac1{L_{\mathrm{eff},i}}
 \sum_\ell M_{i\ell}\min\!\left(\rho_{i\ell}\widetilde A_i,
 \operatorname{clip}(\rho_{i\ell},0.8,1.28)\widetilde A_i\right).
 \label{eq:surrogate}
\end{equation}
Every reported gradient is $\nabla_z\mathcal L$, flattened over completions
and positions, with fixed old log probabilities, masks, and detached weights;
it is not a model-parameter gradient. Token means are computed per completion
before averaging over $I$. The synthetic loss has no KL or entropy regularizer;
adding either is outside the zero-gradient claim. This coupled objective exists
in code, but there is no verified integrated language-model training path.

\section{Trajectory-Spectral Auxiliary Score}

Let $X_i\in\mathbb{R}^{L\times D}$ denote a completion feature sequence.  The
artifact uses either supplied hidden states or a two-channel fixture consisting
of token log probabilities and normalized position.  For
$t_\ell=2\ell/(L-1)-1$ and Legendre polynomials $P_n$, it computes
\begin{equation}
  \widehat c_{i,n}
  = \frac{2n+1}{L_{\mathrm{eff},i}}
    \sum_{\ell=0}^{L-1} M_{i,\ell}X_{i,\ell}P_n(t_\ell),
  \label{eq:discrete-coeff}
\end{equation}
where $M$ is a completion mask and
$L_{\mathrm{eff},i}=\sum_\ell M_{i,\ell}$ is the per-completion active
length.
This is the defective padded-grid rule being audited, not a recommended
projection. Closed endpoint nodes have spacing $2/(L-1)$, inconsistent with
the equal-weight $2/L_{\mathrm{eff},i}$ quadrature. Padding also contracts
the active domain. A prospective active-token trapezoid specification is
\[
\widehat c^{\rm trap}_{i,n}=\frac{2n+1}{L_{\mathrm{eff},i}-1}
\sum_{k=0}^{L_{\mathrm{eff},i}-1}w_kX_{i,\ell(k)}
P_n\!\left(\frac{2k}{L_{\mathrm{eff},i}-1}-1\right),
\]
where $\ell(k)$ enumerates active tokens, endpoint weights are $1/2$ and
interior weights are one. Require $L_{\mathrm{eff},i}\ge\max(2,N)$;
this does not make finite-sample Legendre orthogonality exact.
\textbf{Unmeasured defect.} The masked/padded corruption of
Eq.~\eqref{eq:discrete-coeff} (domain contraction and wrong normalizer
under padding) is asserted analytically here and measured nowhere in this
revision: every replay cell in \S6 uses complete, all-active masks, and
the one unequal-length check validates only the reward-only derivative,
not the spectral projection. The minimum required measurement before any
training claim is masked replay cells (e.g.\ $L_{\mathrm{eff}}\in\{L/4,L/2\}$
per row, including rows of unequal active length in one group) reporting
coefficient error of Eq.~\eqref{eq:discrete-coeff} against the
active-token trapezoid reference above. Until those cells exist, the
defect's magnitude --- as opposed to its existence as a formula mismatch
--- is unknown.
Pairwise distance is
\begin{equation}
  d_N(i,j)=\sum_{n=0}^{N-1}
  \frac{2}{2n+1}\left\|\widehat c_{i,n}-\widehat c_{j,n}\right\|_2^2,
  \qquad
  s_i=\frac{1}{a-1}\sum_{j\in I\setminus\{i\}}d_N(i,j),\quad i\in I.
\end{equation}
Throughout, $d_N$ is a squared weighted modal distance; we never take a square
root.  This matters because rooting is a nonlinear reparameterization that
changes the relative weight of near and far pairs, the within-group variance of
$s_i$, and hence the update induced by any fixed $\lambda$.
\paragraph{The score is outlier-seeking; its sign is a commitment.}
Mean pairwise distance is maximized by the group's most atypical
completion. With $\lambda>0$, an all-wrong group therefore upweights its
most unusual wrong answer, and an all-correct group downweights its most
canonical correct answer relative to atypical correct ones. The first
might be defended as exploration; the second is arguably backwards, and
neither follows from any mechanism linking low-order Legendre modes of a
log-probability trace to correctness. We offer no such mechanism. The
prospective gate in \S7 must therefore test stratum-specific directional
predictions --- positive-$\lambda$ alignment separately in all-correct
and all-wrong strata --- and pre-specify aggregation ablations (median
pairwise distance; signed or centrally referenced variants) so that a
negative gate distinguishes ``trajectory geometry is uninformative'' from
``this outlier-seeking aggregation has the wrong sign.'' The artifact
detaches the combined advantage after standardizing the score; differentiating
through the auxiliary weight would be a different algorithm.

\subsection{What is exact, and what is approximate}

For arbitrary coefficient vectors $c_n$ define the reconstructed function
$\widehat X_N(t)=\sum_{n=0}^{N-1}c_nP_n(t)$.  Orthogonality gives the exact
identity
\begin{equation}
  \int_{-1}^{1}\left\|\widehat X_N(t)\right\|_2^2\,dt
  =\sum_{n=0}^{N-1}\frac{2}{2n+1}\|c_n\|_2^2.
  \label{eq:parseval-truncated}
\end{equation}
Equation~\eqref{eq:parseval-truncated} is exact for the reconstructed expansion.
It does not imply that Eq.~\eqref{eq:discrete-coeff} equals the continuous
projection $c_n=(2n+1)\int_{-1}^1X(t)P_n(t)\,dt/2$ of the unknown
trajectory function. The implemented equally spaced sum is neither
Gauss--Legendre quadrature nor an exact discrete orthogonal transform.
On a fully active grid, its endpoint weighting and $L$ rather than $L-1$
normalizer give first-order error for generic smooth functions (the
mismatched endpoint weight contributes an $O(1/L)$ term that does not
cancel), whereas composite trapezoid quadrature has second-order error
under the usual bounded second-derivative condition. Neither order guarantees accuracy for arbitrary
hidden-state sequences. Under padding, the active domain additionally changes.

The former numerical discretization table is withdrawn because its raw audit
rows and test function are unavailable. The rank limitation remains analytic:
a basis evaluated on $L_{\mathrm{eff}}$ active nodes has rank at most
$L_{\mathrm{eff}}$, so $N>L_{\mathrm{eff}}$ cannot yield full modal rank.
The current implementation does not enforce that admissibility condition
or implement the active-token trapezoid formula above. These remain repairs
required before a training experiment, not properties certified by the replay.

\section{Adaptive Givens Coordinate Elimination}

For feature coordinates $u\ne v$, let $R(u,v,\theta)$ be a planar rotation. With
$\theta=\operatorname{atan2}(x_v,x_u)$ and the sign convention used by the
artifact,
\begin{equation}
  (Rx)_u=\sqrt{x_u^2+x_v^2},
  \qquad
  (Rx)_v=0.
\end{equation}

\begin{lemma}[Norm preservation of the implemented elimination]
A Givens rotation preserves $\ell_2$ norm.  If its angle is selected as above,
projecting coordinate $v$ to zero after the rotation leaves the norm unchanged
in exact arithmetic because that coordinate is already zero.
\end{lemma}

\begin{proof}
Orthogonality gives $\|Rx\|_2=\|x\|_2$. The selected angle annihilates the
$v$th coordinate, so the projection subtracts zero squared magnitude. The
same argument applies inductively only if no later rotation touches a
previously annihilated coordinate before the final projection. This
non-interference condition is required, not implied by orthogonality alone.
\end{proof}

This lemma is narrower than a general claim about ``noise removal.''  The
operation preserves information from the selected coordinate pair by folding
its norm into another coordinate, but it does not establish that the pair is
noise, preserve direction, preserve attention logits, or improve task reward.
Moreover, the hard projection is algebraically redundant: the adaptive
rotation has already made every projected coordinate zero. Floating-point
residue is possible, but the unavailable historical norm-loss sweep cannot
establish its magnitude here. In particular, relative removed energy
$\|(I-\Pi)z\|_2^2/\|z\|_2^2$ and relative norm decrease
$1-\|\Pi z\|_2/\|z\|_2$ for $z=Rx\ne0$ are different metrics.
We retain the exact-arithmetic lemma and withdraw the unsupported numerical
precision, truncation, and invalid-width comparisons. Norm preservation alone
is not evidence of denoising or useful information preservation.

\section{Artifact and Synthetic Stress Tests}

The artifact includes modules for the spectral projection, pairwise score,
Givens elimination, and the modified objective. The entropy gate audited
in this section is vacuous as implemented (0/18 activations; see below)
and is retained only as an audit subject, not as a live component.
Sixteen focused tests pass on
CPU.  They check recurrence values, tensor shapes, nonnegative distance,
rotation invariance, coordinate elimination, fixture construction, output
schema, and a deterministic equal-reward stress test.

\paragraph{Current source versus historical audit.}
The current \path{zvf-program/flagship/pilot/objective.py} explicitly
detaches $A^{\rm reward}+\lambda A^{\rm aux}$ and inserts it in a clipped
score-function loss. It uses sample standard deviation on selected rows,
auxiliary epsilon $10^{-4}$, and a $10^{-12}$ floor. Thus an inference
that the current synthetic harness omits coupling is false; this is still
not evidence of an integrated language-model training experiment. The current
elimination module rejects $N_{\rm noise}\notin[1,D-1]$.
The synthetic objective's gate is per completion:
$p_{i\ell}=\operatorname{softmax}(z_i)_\ell$,
$H_i=-\sum_\ell p_{i\ell}\log(p_{i\ell}+10^{-12})$, and it applies
elimination to each token's feature vector before projection when $H_i<\tau$.
This softmax runs across all padded token positions, without masking or
length normalization, and is not a vocabulary probability distribution.
The gate and rotations affect only auxiliary feature construction; the
surrogate in Eq.~\eqref{eq:surrogate} retains the original $z$ tensor.
This differs from the standalone module's batch-mean gate, and neither
establishes attention-head gating in a trained transformer. For zero-based
rotation step $k=0,\ldots,N_{\rm noise}-1$, $v_k=D-1-k$ and
$u_k=\max(0,v_k-N_{\rm noise})$. Angles use the vector after preceding
rotations, and the annihilated tail is projected only after all steps.
All later coordinate indices are smaller than $v_k$, which ensures the
non-interference condition for valid widths without requiring disjoint pairs.

The replacement replay uses the repository fixture generator with
$G\in\{4,8,16\}$, $L\in\{512,1024,2048\}$, CPU float64, seed 42,
$N=8$, $D=2$, $N_{\rm noise}=1$, and all rows and tokens active.
It sets $z^{\rm old}=0$ and $z_{i\ell}=b_iq_\ell+0.01\xi_{i\ell}$,
where $b=\operatorname{linspace}(-0.1,0.1,G)$,
$q=\operatorname{linspace}(0.1,1,L)$, and $\xi$ is standard Gaussian
noise from a CPU generator reset to seed 42 for each cell. Features are
$(z_{i\ell},t_\ell)$. These are synthetic surrogate inputs, not log
probabilities emitted by a normalized language model.
\textbf{Circularity, stated plainly.} The per-completion offsets $b_i$
are the only systematic between-completion variation, and the second
feature channel $t_\ell$ is identical across completions, so it
contributes exactly zero to every pairwise distance $d_N(i,j)$: the
nominal $D=2$ fixture is effectively $D=1$ for pairwise purposes, and
$s_i$ is, up to noise, a monotone function of the injected offsets. A
nonzero homogeneous gradient and its $\lambda$-linearity then follow
from Eq.~\eqref{eq:combined-advantage} by construction. Any injective
per-completion statistic (e.g.\ token-mean log probability) would pass
the same battery, and no placebo statistic is run in this replay. With
$z^{\rm old}=0$ and $|z|\lesssim0.1$, importance ratios also remain
inside the clip band $[0.8,1.28]$, so the clipping branch of
Eq.~\eqref{eq:surrogate} is never exercised. The replay is therefore an
algebraic identity / code-path check with no failure mode as validation;
it must not be read as evidence that spectral projection recovers task
signal. Adversarial fixtures the score should \emph{not} separate
(identical trajectories at different lengths/padding, permuted-token
controls, ratios driven outside the clip band, a non-spectral placebo
statistic) are required future checks and are not present here. Control rewards alternate
one and zero; homogeneous rewards are all one. The sweep in
Table~\ref{tab:lambda} fixes $G=4,L=512$ and covers every decade from
$10^{-6}$ to $1$ plus zero. It replaces, rather than reproduces, the missing
historical sweep.

\begin{table}[t]
\centering
\caption{New deterministic CPU replay at $G=4,L=512$. Gradient norms use
$\nabla_z\mathcal L$ from Eq.~\eqref{eq:surrogate}; cosines compare the coupled
and reward-only gradients on the alternating-reward control. Values are rounded;
neither quantity is a learning result.}
\label{tab:lambda}
\small
\begin{tabular}{lcc}
\toprule
$\lambda$ & Homogeneous $\|\nabla_z\mathcal L\|_2$ & Cosine vs. reward gradient \\
\midrule
$0$       & $0$                    & $1.000000$ \\
$10^{-6}$ & $1.87848\times10^{-8}$ & $1.000000$ \\
$10^{-5}$ & $1.87848\times10^{-7}$ & $1.000000$ \\
$10^{-4}$ & $1.87848\times10^{-6}$ & $1.000000$ \\
$10^{-3}$ & $1.87848\times10^{-5}$ & $1.000000$ \\
$10^{-2}$ & $1.87848\times10^{-4}$ & $0.999952$ \\
$10^{-1}$ & $1.87848\times10^{-3}$ & $0.995196$ \\
$1$       & $1.87848\times10^{-2}$ & $0.691000$ \\
\bottomrule
\end{tabular}
\end{table}

In this unclipped regime the homogeneous gradient is exactly
$\lambda$ times its $\lambda=1$ value up to floating-point rounding
(norm $0.0187848\lambda$), as differentiating
Eq.~\eqref{eq:combined-advantage} predicts for $\lambda\ge0$ with
detached weights; the replay's role is only to confirm the code
implements that identity.
This proportionality is objective scaling, not recovered task information.
The earlier forward-transform cosine claims are withdrawn with their missing
raw audit; they are not measurements of the present coupled objective.

At the default $\lambda=0.1$, all nine cells have entropy between
$6.23793$ and $7.62457$ nats, above the default threshold $\tau=0.5$.
No row activates the Givens gate, and gated versus ungated gradients are
identical in all 18 reward-regime comparisons. As a branch-execution positive
control, setting $\tau=10$ activates every row and changes the gradient in all
18 comparisons. This demonstrates an executable branch, not a useful denoiser.
Because $H_i$ is the entropy of an $L$-position softmax it scales towards
$\log L$ ($\approx6.24$ at $L=512$), so $\tau=0.5$ is more than an order
of magnitude below every attained value: the 0/18 null is a threshold
scale mismatch (an unnormalized statistic gated by a fixed constant), not
evidence about entropic gating as an idea. As implemented the gate is
vacuous; a mask-aware normalized $H/\log L_{\mathrm{eff}}$ redefinition
is specified prospectively in \S7 and has not been run.
An independent reward-only derivative calculation with six selected rows out
of eight and unequal masked lengths agrees to maximum absolute error
$3.47\times10^{-18}$. End-to-end training cost remains unmeasured.

\paragraph{What the artifact does not show.}
It contains no language-model fine-tuning with this objective, no held-out
GSM8K \citep{cobbe2021gsm8k} or MATH \citep{hendrycks2021math} predictions, no multi-seed accuracy comparison, no calibrated
process target, and no evidence that the auxiliary gradient reduces regret.
Synthetic gradient existence is therefore the endpoint of the present
empirical claim.

\section{Prospective Evaluation}

A publishable mitigation claim requires a prospective experiment in which the
trajectory score is defined before outcomes are observed.

\paragraph{Primary hypotheses.}
\textbf{H1 (alignment):} within homogeneous-reward groups, the preregistered
score predicts a later independently measured outcome, such as
verifier-corrected completion-level success from an independent verifier.
\textbf{H2 (learning):} adding the score improves fixed held-out accuracy or
sample efficiency at bounded KL divergence and equal rollout/token budget.
\textbf{H3 (gating, conditional and currently vacuous):} only after the
gate is redefined in normalized mask-aware form ($H/\log
L_{\mathrm{eff}}$) and shown to fire at a calibration-fixed rate would
``the Givens component adds benefit beyond the same spectral auxiliary
without gating'' be a testable hypothesis. For the implementation
audited here H3 is demoted: the component never fires, so there is
nothing to ablate.

\paragraph{Minimum design.}
Use one canonical, completion-masked GRPO implementation with clipping and
reference KL; one public model family at $1$--$2$B and $7$--$8$B scales; GSM8K
and a harder MATH subset \citep{cobbe2021gsm8k,hendrycks2021math}; and a contamination control.  At both scales and on
both tasks compare: unmodified GRPO; spectral auxiliary; Givens-only; combined;
a random score matched in within-group distribution and variance; an
entropy-shaped advantage; a process/verifier-repair signal where available;
and the simple alternatives of dynamic resampling and larger $G$ at equal
tokens.  Give the baseline back the method's extra compute as a budget-matched
control.  Freeze the sampler, reward parser, optimizer, prompts, maximum
length, group size, evaluation set, checkpoint rule, $\lambda$ schedule, and
token budget before inspecting outcomes. Two cheaper branches get equal
billing in the decision tree, since two of the four latent causes of
homogeneity are verifier defects: verifier repair and prompt-difficulty
curation, evaluated on the same frozen rollouts (the alignment labels
below double as a verifier-error audit), must be compared before any
objective-side remedy is credited. Opportunity cost is likewise explicit:
this program is worth its tokens only if its alignment effect beats
spending the same budget on larger $G$ or dynamic resampling
\citep{yu2025dapo}; projection overhead at realistic $(L,D)$ is
currently unmeasured and must be reported before any training comparison.
Before any fine-tuning design, an offline hidden-state trace study on
frozen checkpoints (no training) must establish whether the score's
behavior on the unexercised hidden-states path resembles the fixture at
all.

Run the cheapest decisive gate first.  From early, middle, and late frozen
checkpoints at each scale, collect at least $2{,}000$ all-correct and $2{,}000$
all-wrong homogeneous groups.  Without using $s_i$ to construct labels, obtain
completion-level verifier-corrected success. Groups with constant secondary
labels cannot identify within-group ranking; report their frequency and do
not assign them an AUC. Fit a length/entropy residualization on separate
calibration prompts, then evaluate within-group concordance on held-out
prompts. The $2{,}000+2{,}000$ group counts above are placeholders, not
a justified size: before freezing, size the gate by target
prompt-clustered interval width, and require a pre-specified minimum
fraction of discordant (non-constant secondary label) groups per stratum,
so the gate cannot pass on a small unrepresentative subset. Require the lower bound of a prompt-cluster bootstrap $95\%$
interval to exceed $0.5$, matching positive $\lambda$, separately for
all-correct and all-wrong strata at both scales and tasks. Prompt-level
pass@$k$ on independent rollouts is not a completion-level alignment target.
If this gate fails, do not train the auxiliary objective; report
the negative alignment result.

If it passes, screen with five paired seeds and expand using
\begin{equation}
 n\ \geq\ \frac{(z_{1-\alpha/2}+z_{1-\beta})^2s_d^2}{\delta^2},
 \label{eq:power}
\end{equation}
where $z_q$ is the standard normal $q$-quantile, $\alpha$ is the two-sided
significance level, $1-\beta$ is the target power,
$s_d$ is the standard deviation of within-seed differences, and $\delta$
is the preregistered smallest effect worth claiming.  With two-sided
$\alpha=0.05$, Eq.~\eqref{eq:power} is only a normal-approximation planning
heuristic, not an exact five-seed power guarantee. Exact noncentral paired-$t$
power with multiplicity adjustment must be computed before a confirmatory
sample size is frozen. The four proposed primary contrasts are combined
versus unmodified GRPO, budget-matched GRPO, variance-matched random score,
and spectral-auxiliary-only; use Holm adjustment across them. This is a
prospective design specification, not evidence of a registered executed trial.
Report per-seed learning curves, fixed-budget
held-out accuracy, pass@$k$, score calibration, reference KL, clip fraction,
output length, GPU-hours, tokens, peak memory, all-correct/all-wrong group
rates, and every failed or collapsed run.

Before that comparison, retain the existing width guard and explicit auxiliary
coupling, but repair the padded spectral grid using active-token coordinates
or a discrete orthogonal basis. Make the per-completion gate mask-aware and
normalize its entropy by $\log L_{\mathrm{eff}}$ for lengths greater than one,
with an explicit singleton policy. Fix its threshold on calibration data.
Integrate and verify this contract in a language-model trainer; neither
attention-head gating nor end-to-end loss integration is established by the
synthetic replay. These are prospective changes, not silently applied repairs
to the modules whose current behavior is reported above.

\paragraph{Falsification criteria.}
Reject the mitigation claim if the score fails out-of-sample alignment, if a
variance-matched random score performs similarly, if gains vanish under paired
seeds, equal tokens, bounded KL, or the budget-matched baseline, if a simple
resampling/larger-group control matches it, or if the Givens ablation adds no
benefit.  A nonzero gradient, norm identity, or passing synthetic test is never
a success criterion.

\section{Limitations and Responsible Claiming}

The auxiliary score is self-referential because it is computed from policy
features or log probabilities.  It may reward novelty, length, formatting, or
instability rather than correctness, and as defined it is outlier-seeking:
with positive $\lambda$ it penalizes the most canonical member of an
all-correct group (\S4). The masked/padded defect of
Eq.~\eqref{eq:discrete-coeff} is unmeasured by the all-active replay, the
entropy gate is vacuous as implemented, and the $\lambda$-sweep is an
identity check on a fixture that injects the recovered signal (\S6).
The within-group framing also discards prompt-level information:
homogeneous groups can still signal prompt difficulty for curriculum or
filtering even when they carry no within-group contrast, a use this audit
does not evaluate. Under prompt-dependent difficulty the binding
constraint may be the prompt distribution rather than advantage algebra.  Hidden-state distances are model- and
layer-dependent.  The equispaced Legendre approximation is sensitive to masks
and sequence length.  The adaptive Givens transform is data dependent and may
change directions even when it preserves norms.  Finally, the current stress
tests use idealized tensors and cannot estimate end-to-end memory, latency, or
learning effects.

\section{Conclusion}

Homogeneous rewards remove the centered reward signal.  A trajectory-derived
score can replace that absence with an auxiliary preference, but doing so
changes the objective and creates a new validation burden. The current
synthetic objective is not a verified integrated trainer. The paper provides
an exact continuous identity for the
reconstructed expansion, a norm-preserving adaptive coordinate elimination,
and deterministic synthetic stress tests. There is no evidence here of
better policy learning.  By separating those statements and specifying
falsification criteria, this revision turns an unsupported performance story
into a documented negative result plus a testable methods proposal. That
is the whole claim: a prototype autopsy for a negative-results /
reproducibility audience, whose reusable parts are the credited floor
check, the exactness boundary, and the claim-to-file discipline --- not
the spectral score, which remains unvalidated, outlier-seeking, and
untested under masking, gating, or training.

\appendix
\section{Claim-to-File Index}
\label{app:artifact}

Repository: \url{https://github.com/arvindcr4/tinker-rl-lab}.
All repository paths below are relative to its root; \path{pilot/} abbreviates
\path{zvf-program/flagship/pilot/}. The replay is saved as
\path{platform_hybrid/experiments/results/review_2026-10-18/Spectral_Entropy/replay.json}.
Its \path{source_sha256} map identifies the four implementation files used.
The replay was regenerated from the committed tree with the command below
(git commit \texttt{39f7396c6043ee83e9584bb67cfce25f1d0efc08} at
regeneration); every numerical field (grid, $\lambda$-sweep including
Table~\ref{tab:lambda}, active-row error) is identical before and after
regeneration --- only the previously stale \path{objective.py} hash
changed, and all four hashes now verify against the committed sources.
Building this file additionally requires \path{neurips_2026.sty}, which
currently lives under \path{zvf-program/flagship/paper/} and is not
vendored here (author TODO at submission, with anonymization, once the
workshop venue is bound; see file header).
\path{python}, \path{torch}, and \path{dtype} record the runtime (Python
3.14.7, PyTorch 2.9.0, float64 CPU). Reproduction requires these sources and a
Python environment with PyTorch, not access to the missing historical audit.
No DOI or independently archived release is claimed. The historical bundle's
case count, checksum, numerical tables, and purported internal-consistency
verdict are not certified here and are no longer evidence for this revision.

From the repository root, run:
\begin{quote}\small\ttfamily
python3 tools/review\_ledger/Spectral\_Entropy\_recompute.py\\
PYTHONPATH=zvf-program/flagship python3 -m unittest\\
\hspace*{1em}pilot.test\_spectral\_attention pilot.test\_entropic\_gating\\
\hspace*{1em}pilot.test\_benchmark\_spectral\_harness
\end{quote}
Join the three displayed unittest lines into one shell command.
The numerical ledger in
\path{reports/deepreview/ledgers/Spectral_Entropy.tsv} records both recomputed
claims and withdrawals. Table~\ref{tab:claim-index} gives field-level pointers
for the retained empirical results; analytical identities and future design
choices are not measurements.

\begin{table}[h]
\centering
\caption{Traceability index for the revision's empirical quantities.}
\label{tab:claim-index}
\scriptsize
\begin{tabularx}{\textwidth}{p{.27\textwidth}X}
\toprule
Claim & Artifact pointer \\
\midrule
16 focused repository tests &
\path{pilot/test_spectral_attention.py},
\path{pilot/test_entropic_gating.py}, and
\path{pilot/test_benchmark_spectral_harness.py} \\
Coupling sweep &
\path{replay.json:lambda_sweep_G4_L512[]} fields \path{lambda},
\path{homogeneous_gradient_norm}, \path{heterogeneous_cosine_to_baseline} \\
Entropy range and inactive default gate &
\path{replay.json:grid[]} fields \path{entropy_min}, \path{entropy_max},
\path{default_gate_count}; for each of \path{alternating}, \path{homogeneous},
\path{default_gradient_max_difference} \\
Forced-active positive control &
\path{replay.json:grid[]} under each reward regime:
\path{forced_gate_count}, \path{forced_gradient_max_difference};
threshold in \path{forced_control_tau} \\
Unequal-length active-row derivative &
\path{replay.json:active_row_reference_max_error}; independent formula in
\path{tools/review_ledger/Spectral_Entropy_recompute.py} \\
\bottomrule
\end{tabularx}
\end{table}

\bibliographystyle{plainnat}
\bibliography{spectral_audit}

\end{document}
